\scanpage{138}
\begin{center}- \mspage{137} -\end{center}\margin{IX}

sur $K$, qui jouera dans le cas présent le même rôle que 11.2 dans le cas traité précédemment :

\uline{Théorème} 11.5. \uline{Supposons} $S$ \uline{hensélien}, \uline{et que} $A_K$ \uline{ait réduction semi-stable sur} $S$. \uline{Soit} $\ell$ \uline{un nombre premier}, \uline{et considérons l'accouplement de monodromie} (9.1.2), \uline{d'où un homomorphisme canonique}
\[
 \widetilde u_\ell:\underline M_\ell\longrightarrow\check{\underline M}'_\ell.
\]
\uline{Soit} $\Phi(\ell)=\operatorname{coker}\widetilde u_\ell$ ; \uline{alors on a un isomorphisme canonique}
\[
(11.5.1)\qquad\Phi_{\unsure{o|0}}(\ell)\simeq\Phi(\ell)\times_Ss.
\]

De cet énoncé résulte immédiatement la définition d'un accouplement parfait (11.4.1), puisque on trouve, en appliquant 11.5 à $A'_K$ au lieu de $A_K$ et utilisant (10.2.2), qu'on a aussi un isomorphisme canonique
\[
\Phi'_{\unsure{o|0}}(\ell)\simeq\Phi'(\ell)\times_Ss,\quad\text{où}\quad
\Phi'(\ell)=\operatorname{coker}({}^t\widetilde u_\ell:\underline M'_\ell\longrightarrow\check{\underline M}_\ell).
\]

\uline{Remarques} 11.5.2. a) Le théorème 11.5 implique notamment que $\widetilde u_\ell$ est une isogénie i.e. que $u_\ell$ est séparant, y compris dans le cas $\ell=p$, où cette assertion n'est pas triviale sur la définition de 9.6. On notera que nous n'utilisons pas ici le théorème (nettement plus profond) 10.4, qui sera prouvé au par. suivant.

b) Utilisant 10.4, qui nous fournit un accouplement de monodromie $u:\underline M\otimes\underline M'\longrightarrow\mathbb Z_S$, et désignant par $\Phi$ le conoyau de l'homomorphisme correspondant
\[
 \widetilde u:\underline M\longrightarrow\check{\underline M}',
\]
\figure{\begin{tikzpicture}[scale=.4]
\draw (0,0)--(.2,.1)--(.35,-.1)--(.5,.12)--(.7,-.08)--(.9,.1);
\draw (2,0)--(3,0);\draw (2.2,.12)--(2,0)--(2.2,-.1);
\draw (3.4,-.1)..controls(3.15,-.35)and(3.45,-.4)..(3.55,-.22);
\draw (1,-1.8)--(1.14,-1.9)--(.9,-2.1)--(1.05,-2.2);
\end{tikzpicture}}
\margin{455}\margin{137}
